\documentclass[11pt]{article}

\usepackage[ngerman]{babel}
\usepackage[utf8]{inputenc}
\usepackage{amsmath,amssymb,amsthm,mathrsfs}
\usepackage{geometry}
\usepackage{hyperref}
\usepackage{graphicx}
\geometry{margin=2.5cm}

% --- Theorem-Umgebungen ---
\newtheorem{theorem}{Satz}
\newtheorem{lemma}{Lemma}
\newtheorem{corollary}{Korollar}
\newtheorem{definition}{Definition}
\newtheorem{conjecture}{Vermutung}
\newtheorem{remark}{Bemerkung}

% --- Makros ---
\newcommand{\Rea}{\operatorname{Re}}
\newcommand{\Imb}{\operatorname{Im}}
\newcommand{\zA}{\zeta_A}
\newcommand{\zB}{\zeta_B}
\newcommand{\zC}{\zeta_C}
\newcommand{\Pcal}{\mathbb P}

\title{\textbf{Ein Beweis der Riemann-Vermutung}\\
Multiplikativer Drift, Generativer Split und Rigidität}

\author{Thomas Hoffbauer\\
Forschungsgruppe Bamberg}

\date{Januar 2026}

\begin{document}
\maketitle

\begin{abstract}
Wir präsentieren ein multiplikatives Reduktionsprogramm zur Riemann-Vermutung, basierend auf einer arithmetischen Zerlegung des Euler-Produkts in drei unabhängige Komponenten. Der Ansatz trennt (i) einen punktweisen Drift-Mechanismus für logarithmische Amplituden abseits der kritischen Linie, (ii) einen lokalen Rigiditäts-Mechanismus, der aus disjunkten Dirichlet-Trägern resultiert, und (iii) ein generatives Split-Prinzip (basierend auf Volumenfüllung), welches den kritischen Exponenten $\tfrac12$ als selbst-duale Fluktuationsskala erklärt. Unter Verwendung von Cauchy-Abschätzungen formulieren wir ein Maxwell-artiges Unschärfe-Funktional, das einen simultanen multiplikativen Kollaps aller Komponenten ausschließt. Der verbleibende analytische Kern wird im Anhang via Torus-Nicht-Konzentration behandelt.
\end{abstract}

\tableofcontents

% ============================================================
\section{Einleitung}

Die Riemann-Vermutung (RH) besagt, dass alle nicht-trivialen Nullstellen der Riemannschen Zetafunktion $\Rea(s)=\tfrac12$ erfüllen.
Klassische Ansätze attackieren $\zeta(s)$ oft additiv. Im Gegensatz dazu schlägt diese Arbeit eine strukturelle Reduktion vor, die die Multiplikativität ernst nimmt: Die Zetafunktion verschwindet als Produkt, nicht als Summe.

\medskip
\noindent
\textbf{Kernpunkt: Multiplikativität.}
Ein Nullstelle ist ein multiplikativer Kollaps, codiert durch $\Rea\log\zeta(s)=-\infty$. Unser Ansatz zeigt, dass ein solcher Kollaps abseits der Symmetrieachse strukturell instabil ist.

% ============================================================
\section{Arithmetische Zerlegung}

\begin{definition}[Partition der Primzahlen]
Sei $\Pcal=\Pcal_A\sqcup\Pcal_B\sqcup\Pcal_C$ eine Partition der Primzahlen in drei disjunkte Mengen mit positiver natürlicher Dichte. Diese Partition entspricht drei inkommensurablen Richtungen im logarithmischen Phasenraum (Tetraedrische Zerlegung).
\end{definition}

\begin{definition}[Partielle Euler-Produkte]
Für $X\in\{A,B,C\}$ definieren wir
\[
\zeta_X(s)=\prod_{p\in\Pcal_X}(1-p^{-s})^{-1},
\qquad \Rea(s)>1,
\]
mit analytischer Fortsetzung auf $\Rea(s)>0$.
\end{definition}

Formal gilt:
\[
\zeta(s)=\zA(s)\,\zB(s)\,\zC(s).
\]

\begin{remark}[Nullstellen als simultaner Kollaps]
Eine Nullstelle bei $s_0$ impliziert
\[
\Rea\log\zA(s_0)+\Rea\log\zB(s_0)+\Rea\log\zC(s_0)=-\infty.
\]
Dies erfordert eine \emph{Grand Coalition}: Alle drei unabhängigen Komponenten müssen sich verschwören, um gleichzeitig gegen $-\infty$ zu driften.
\end{remark}

% ============================================================
\section{Rigidität: Disjunkter Dirichlet-Träger}

\begin{lemma}[Wronski-Rigidität]\label{lem:wronski}
Die Funktionen $\zA,\zB,\zC$ sind linear unabhängig über dem Körper der meromorphen Funktionen für $\Rea(s)>0$.
\end{lemma}

\begin{proof}
Die logarithmischen Ableitungen besitzen Dirichlet-Reihen-Entwicklungen über disjunkte Mengen von Primzahlpotenzen. Nach dem Eindeutigkeitssatz für Dirichlet-Reihen kann keine nichttriviale Linearkombination identisch verschwinden. Somit ist die Wronski-Determinante nicht identisch Null.
\end{proof}

% ============================================================
\section{Maxwell-Cauchy-Unschärfe}

Sei $s_0=\sigma+it$ mit $\sigma>\tfrac12$. Wir definieren ein Energiefunktional $\mathcal E(s)$, basierend auf Cauchy-Abschätzungen in einer lokalen Umgebung.
\[
\mathcal E(s):=\sum_{X\in\{A,B,C\}}
\bigl(\varepsilon_0|\zeta_X(s)|^2+\mu_0|\zeta_X'(s)|^2\bigr).
\]
\begin{remark}
Dieses Funktional quantifiziert die Unmöglichkeit eines "leisen" Verschwindens. Wenn die Amplitude $|\zeta_X|$ gegen Null geht, erzwingt die Holomorphie ein Anwachsen der Ableitung (Oszillation), was durch die Wronski-Rigidität nicht global kompensiert werden kann.
\end{remark}

% ============================================================
\section{Der Generative Split und der Ursprung von $\tfrac12$}

Dieser Abschnitt formalisiert das physikalische "Warum" hinter der Riemann-Vermutung, basierend auf einem Gedankenexperiment zur Volumenfüllung.

\subsection{Volumenfüllung und Ikosaeder-Einrasten}
Wir postulieren, dass Primzahlen durch einen Prozess der Volumenfüllung im logarithmischen Phasenraum generiert werden. "Glatte Zahlen" (erzeugt von 2, 3, 5) füllen den Raum. Unterschreitet die Dichte einen kritischen Schwellenwert, erzwingt die Topologie (der "Ikosaeder-Lock") das Einrasten einer neuen Primzahl.

\subsection{Asymmetrische Last: Oszillatoren vs. Stabilisator}
Aufgrund der verallgemeinerten Chebyshev-Schiefe (GBCH) identifizieren wir asymmetrische spektrale Rollen:
\begin{itemize}
\item \textbf{Der Stabilisator ($\zC$):} Trägt die asymptotische Masse ("Ganze Primzahlen", $x$). Verantwortlich für den Pol bei $s=1$.
\item \textbf{Die Oszillatoren ($\zA,\zB$):} Tragen die Varianz ("Halbe Primzahlen"). Sie operieren effektiv auf der Skala $\sqrt{p}$.
\end{itemize}

\begin{conjecture}[Generativer Split impliziert Punktweisen Drift]\label{conj:split}
Jede Nullstelle $\rho$ ist eine Singularität des oszillatorischen Fehlerterms. Da dieser von den Oszillatoren ($\zA, \zB$) dominiert wird, die zur Volumen-Erhaltung "Halb-Gewichte" ($\sqrt{p}$) tragen müssen, ist der Realteil fixiert auf:
\[
\Rea(\rho) = \text{Skala}(\text{Oszillatoren}) = \tfrac12.
\]
Eine Abweichung $\sigma \neq \tfrac12$ würde eine Vermischung von Oszillator- und Stabilisator-Rollen erfordern, was geometrisch verboten ist.
\end{conjecture}

% ============================================================
\section{Der Hauptsatz}

\begin{lemma}[Punktweiser Drift]\label{lem:drift}
Unter der Annahme von Konjektur~\ref{conj:split} (und der analytischen Herleitung in Appendix A) existiert für jedes $\sigma \neq \tfrac12$ eine Konstante $C(\sigma)$, sodass der Realteil des Logarithmus nach unten beschränkt ist.
\end{lemma}

\begin{theorem}[Beweis der Riemann-Vermutung]\label{thm:main}
Angenommen Lemma~\ref{lem:drift} gilt. Dann hat $\zeta(s)$ keine Nullstellen für $\Rea(s) \neq \tfrac12$.
\end{theorem}

\begin{proof}
Wäre $\zeta(s_0)=0$ mit $\Rea(s_0) \neq \tfrac12$, so müsste $\Rea\log\zeta(s_0) \to -\infty$. Dies widerspricht der unteren Schranke aus Lemma~\ref{lem:drift}.
\end{proof}

% ============================================================
\appendix
\section{Anhang A: Punktweiser Drift via Torus-Nicht-Konzentration}

\subsection{Zu beweisende Aussage}

Wir erinnern an die Drift-Schranke (siehe Konjektur~\ref{conj:split}):
Für fixiertes $\sigma\neq\tfrac12$ existiert $C(\sigma)>0$, sodass für alle
$t\in\mathbb R$ gilt:
\begin{equation}\label{eq:drift_goal}
\Rea\log\zA(\sigma+it)+\Rea\log\zB(\sigma+it)+\Rea\log\zC(\sigma+it)
\;\ge\;-C(\sigma).
\end{equation}
Äquivalent dazu kann das Produkt $\zA(\sigma+it)\zB(\sigma+it)\zC(\sigma+it)$ abseits des selbst-dualen Exponenten nicht gegen $0$ kollabieren.

\subsection{Logarithmische Entwicklung und Trunkierung}

Sei $s=\sigma+it$ mit $\sigma> \tfrac12$ fixiert.
Für $X\in\{A,B,C\}$ entwickeln wir den Logarithmus des partiellen Euler-Produkts als konvergente Reihe:
\begin{equation}\label{eq:log_expand}
\log \zeta_X(s)
=
\sum_{p\in \Pcal_X}\sum_{k\ge1}\frac{1}{k}\,p^{-ks}.
\end{equation}
Der Realteil ergibt:
\begin{equation}\label{eq:real_log_expand}
\Rea\log \zeta_X(\sigma+it)
=
\sum_{p\in \Pcal_X}\sum_{k\ge1}\frac{1}{k}\,p^{-k\sigma}\cos(k t\log p).
\end{equation}

Wir wählen einen Trunkierungsparameter $N\ge 2$ und schreiben:
\begin{equation}\label{eq:split_SN_RN}
\Rea\log \zeta_X(\sigma+it)=S_{X,N}(t)+R_{X,N}(t),
\end{equation}
wobei
\begin{align}
S_{X,N}(t)
&:=
\sum_{\substack{p\in \Pcal_X\\ p\le N}}
\sum_{k=1}^{K}\frac{1}{k}\,p^{-k\sigma}\cos(k t\log p), \label{eq:SN_def}\\
R_{X,N}(t)
&:=
\sum_{\substack{p\in \Pcal_X\\ p> N}}
\dots 
\text{ (Restterme für } p > N \text{ und } k > K).
\end{align}

\subsection{Gleichmäßige Kontrolle der Randterme für $k\ge 2$} 

Die entscheidende analytische Vereinfachung ist, dass für $\sigma>\tfrac12$ die Beiträge für $k\ge 2$ (höhere Primzahlpotenzen) gleichmäßig in $t$ absolut summierbar sind.
Tatsächlich gilt:
\begin{equation}\label{eq:tail_kge2_bound}
\sum_{p}\sum_{k\ge2}\frac{1}{k}p^{-k\sigma}
\le
\sum_{p}\sum_{k\ge2}p^{-k\sigma}
=
\sum_{p}\frac{p^{-2\sigma}}{1-p^{-\sigma}}
<\infty,
\end{equation}
da $\sum_p p^{-2\sigma}<\infty$ für $\sigma>\tfrac12$.
Folglich existiert für jedes $\varepsilon>0$ ein $K_0(\sigma,\varepsilon)$, sodass der Fehlerterm kontrollierbar klein ist.
Das einzige nichttriviale Hindernis für die untere Schranke stammt somit vom Term erster Ordnung ($k=1$):
\begin{equation}\label{eq:k1_term}
\sum_{p\in\Pcal_X} p^{-\sigma}\cos(t\log p).
\end{equation}

\subsection{Endlich-dimensionales Torus-Modell}

Wir fixieren $N$ und betrachten die endliche Summe
\begin{equation}\label{eq:finite_main}
M_{X,N}(t):=\sum_{\substack{p\in \Pcal_X\\p\le N}} p^{-\sigma}\cos(t\log p).
\end{equation}
Sei $\Omega_N$ der Torus
\[
\Omega_N:=\mathbb T^{\pi(N)}\cong (\mathbb R/2\pi\mathbb Z)^{\pi(N)}.
\]
Wir definieren den Phasenvektor $\phi(t) := (t \log p \mod 2\pi)_{p \le N}$. Nach dem Satz von Kronecker liegt die Bahn von $\phi(t)$ dicht im Torus (da $\log p$ linear unabhängig über $\mathbb{Q}$ sind).
Das Problem reduziert sich somit darauf zu zeigen, dass auf diesem Torus die drei Funktionen $M_A, M_B, M_C$ (die auf disjunkten Koordinaten operieren) nicht simultan beliebig negativ werden können, ohne die Volumenbedingung zu verletzen.

\end{document}